\chapter{Varicose Vein Ligation and Stripping}

\section*{Introduction \& Clinical Indications}
Ligation and stripping removes an incompetent superficial truncal vein and its refluxing segment. It is one treatment for symptomatic varicose veins, skin changes, superficial vein thrombosis, bleeding, or venous ulceration after duplex assessment. Modern management usually considers endothermal ablation first, foam sclerotherapy when ablation is unsuitable, and surgery when those options are unsuitable or not preferred.

\section*{Relevant Surgical Anatomy}
The great saphenous vein runs medially and joins the common femoral vein at the saphenofemoral junction. The small saphenous vein runs posteriorly and commonly joins the popliteal vein. Perforator veins connect superficial and deep systems. The saphenous nerve is vulnerable near the distal great saphenous vein and the sural nerve near the distal small saphenous vein. Duplex mapping determines the refluxing segments and the safe treatment plan.

\section*{Preoperative Workup \& Preparation}
Duplex ultrasound confirms reflux, maps truncal and tributary anatomy, and excludes relevant deep venous disease. The history includes ulceration, bleeding, thrombosis, arterial disease, pregnancy, prior venous procedures, medications, and mobility. Consent covers endovenous alternatives, bruising, hematoma, infection, nerve injury, deep vein thrombosis, pulmonary embolism, recurrence, and residual tributaries. Interventional treatment is generally deferred during pregnancy except in exceptional circumstances; compression may be used for symptoms. Routine laboratory tests and antibiotics are selective.

\section*{Surgical Technique \& Procedure}
\begin{enumerate}
  \item The patient is positioned supine for the great saphenous vein or prone for some small saphenous procedures. Anesthesia may be local with sedation, regional, or general.
  \item The mapped junction and vein are exposed through small incisions. The truncal vein is ligated at the appropriate junction while preserving the deep vein.
  \item A stripper is passed through the selected segment and the vein is removed, or a limited phlebectomy is performed for tributaries. The extent is based on anatomy and symptoms rather than routine stripping to the ankle.
  \item Hemostasis is confirmed, incisions are closed, and a dressing or short-term compression is applied according to protocol.
\end{enumerate}

\section*{Complications \& Risks}
Risks include bruising, hematoma, wound infection, lymphatic leak, sensory nerve injury, superficial thrombophlebitis, deep vein thrombosis, pulmonary embolism, skin injury, recurrence, and persistent symptoms from untreated or newly incompetent veins.

\section*{Postoperative Care \& Recovery}
Patients are encouraged to walk early and receive wound and analgesia instructions. Compression after treatment is individualized and, when used after intervention, is generally short term rather than indefinite. Review is arranged for wound healing, persistent symptoms, thrombosis concerns, and residual or recurrent veins. New unilateral swelling, chest pain, breathlessness, heavy bleeding, fever, or severe pain requires urgent assessment.

\section*{Outcomes \& Prognosis}
Symptoms often improve after appropriate treatment of reflux, but recurrence and residual tributaries can occur. Outcomes depend on venous anatomy, severity of chronic venous disease, ulcer care, and the selected intervention. Ligation and stripping remains effective when appropriately selected but is not universally first-line therapy.

The indication should be documented in functional terms: aching, heaviness, edema, recurrent superficial thrombophlebitis, bleeding, skin inflammation, or an ulcer associated with superficial reflux. Cosmetic concerns may also be addressed, but the patient should understand that treatment of a truncal vein may not remove every visible tributary and that additional phlebectomy or sclerotherapy may be needed. A normal duplex study should prompt consideration of another cause of symptoms rather than unnecessary vein removal.

The saphenofemoral and saphenopopliteal junctions require careful identification because mistaken ligation of a deep vein can cause major morbidity. Variations in the small saphenous termination, duplicated trunks, accessory saphenous veins, and pelvic or perforator reflux can explain persistent disease. Duplex mapping should be performed with the patient in a position that demonstrates reflux and should record the segments treated, the deep venous system, and any thrombosis.

Patients with chronic venous insufficiency benefit from skin moisturization, treatment of dermatitis, appropriate compression, exercise, and ulcer care before and after intervention. Compression is contraindicated or modified in significant peripheral arterial disease, so pulses and arterial status should be considered. A venous ulcer that is infected, rapidly enlarging, or associated with systemic illness requires urgent wound and medical assessment rather than routine elective stripping.

The recovery plan includes walking, hydration, analgesia, and observation of the wounds and limb. Bruising can track along the vein and may last several weeks. Persistent numbness, burning, swelling, or pain should be reviewed, especially if it follows the distribution of a cutaneous nerve. Duplex ultrasound is obtained when there is concern for deep-vein thrombosis, pulmonary embolism, recurrent reflux, or an unexpected change in symptoms. Long-term recurrence reflects both the treated segment and progression of venous disease; a durable result requires continued management of the underlying venous risk factors.

\section*{Counseling and Long-Term Follow-Up}
The consultation should distinguish venous symptoms from cosmetic concerns and document the patient's baseline function. A venous clinical severity score, ulcer measurements, edema, skin pigmentation, eczema, lipodermatosclerosis, and prior thrombosis can help establish a baseline. Duplex should show the refluxing segment and confirm that the deep veins are patent or identify obstruction requiring a different plan. A patient with severe edema from iliac obstruction may not improve after treating only a superficial vein.

The choice of procedure is shared. Endothermal ablation usually avoids groin dissection and extensive bruising, while foam sclerotherapy can treat selected tributaries or recurrent segments and surgery can address tortuous veins or anatomy unsuitable for a catheter. Risks, recurrence, need for staged treatment, residual cosmetic veins, nerve symptoms, and the possibility that symptoms will persist should be discussed. Treatment of a visible vein is not an indication to remove a normal deep vein or a superficial vein needed for future bypass access without careful consideration.

Bleeding from a varicosity is managed by lying down, elevating the leg, and applying firm direct pressure for a sustained period; emergency services are needed if bleeding is heavy or does not stop. Patients with ulcers should receive a wound plan, infection precautions, and compression advice tailored to arterial status. New unilateral swelling, warmth, calf pain, chest pain, syncope, or breathlessness requires urgent assessment for thromboembolism. Fever, spreading redness, purulent drainage, or severe postoperative pain may indicate infection or another complication.

Long-term results depend on the underlying chronic venous disease. Weight management, exercise of the calf pump, avoidance of prolonged immobility, skin care, and treatment of recurrent reflux support symptom control. Pregnancy and future hormonal changes may produce new veins even after technically successful treatment. Recurrence should be investigated with duplex rather than assumed to be a failed operation. When the treated vein was used as a conduit for bypass or future access, the plan should be documented and alternative access discussed before intervention.

\section*{Patient Education}
Patients should understand that the treated vein is not needed for normal circulation because the deep system carries the principal venous return, but the deep veins must remain patent. Walking is encouraged unless another condition prevents it. Patients should follow the prescribed wound and compression plan and avoid self-treatment of a painful swollen leg without assessment. The appearance of bruising improves gradually and does not by itself indicate treatment failure.

\section*{Patient Selection Summary}
The best candidates are those with symptoms or complications attributable to documented superficial reflux and a treatment plan that addresses the relevant venous segments. The presence of cosmetic veins alone calls for a different discussion from an ulcer or recurrent bleeding. A treatment plan should identify whether the goal is symptom control, ulcer healing, prevention of bleeding, or appearance and should state what further procedures may be needed.

\section*{Documentation and Safety}
The operative note records the duplex-mapped segments, junction treated, deep-vein status, phlebectomy sites, nerve and skin precautions, compression plan, and any complications. Patients should know how to walk safely, care for wounds, recognize thrombosis or infection, and obtain urgent help for chest symptoms or uncontrolled bleeding. Recurrence is assessed with examination and duplex rather than judged from early bruising.

\section*{Practical Follow-Up}
Early walking and wound review are important. Bruising may persist, but increasing unilateral swelling, severe pain, fever, uncontrolled bleeding, chest pain, or breathlessness requires urgent evaluation. Recurrent veins are assessed after healing with examination and duplex imaging when indicated.

\section*{Safety Summary}
Treatment should target duplex-confirmed clinically relevant reflux and should preserve the deep venous system. Endovenous treatment, sclerotherapy, surgery, and conservative care are alternatives whose risks and goals should be discussed individually.

\section*{Final Safety Note}
Treatment targets documented clinically relevant reflux, not every visible vein. Walking, wound care, and individualized compression are part of recovery; new swelling, chest symptoms, fever, or uncontrolled bleeding requires urgent assessment.

\section*{Completion Checklist}
Record duplex-defined reflux, deep-vein status, treatment goal, procedure, compression plan, and warning symptoms. Recurrence or persistent pain should be assessed clinically and with duplex when indicated, not assumed to represent failure.

\section*{Handoff}
The venous handoff should record treated segments, compression advice, walking instructions, and urgent symptoms of thrombosis, infection, or embolism.

\section*{References}

\section*{Additional Clinical Considerations}
Treatment begins with duplex-defined reflux and a discussion of symptoms and goals. Visible veins alone do not always explain leg pain, and other causes such as arterial disease, neuropathy, musculoskeletal pain, lymphedema, and deep venous obstruction should be considered. Compression, exercise, weight management, elevation, skin care, and ulcer treatment may be useful, but compression should not substitute for vascular assessment when there is suspected arterial insufficiency.

Current treatment pathways generally favor endothermal ablation for suitable truncal reflux, with ultrasound-guided foam sclerotherapy or surgery used when anatomy, prior treatment, device availability, patient preference, or risk makes ablation unsuitable. Stripping may be chosen for a large, tortuous, superficial vein or when an endovenous approach cannot safely treat the segment. The goal is to eliminate clinically relevant reflux while preserving useful deep and superficial venous drainage; routine removal of every visible vein is neither necessary nor always safe.

Patients with active venous ulceration need a plan for wound care, infection assessment, compression, and reflux treatment. A bleeding varicosity requires immediate elevation and direct pressure; recurrent bleeding warrants urgent vascular assessment. A painful, red, cord-like vein may represent superficial thrombophlebitis, but unilateral swelling, chest pain, or breathlessness raises concern for deep-vein thrombosis or pulmonary embolism and requires urgent evaluation. Anticoagulation decisions are based on thrombus extent and risk, not on surgery alone.

After stripping, bruising and tightness along the treated track are common. Walking is encouraged when safe, while strenuous activity and wound care follow the procedural plan. Numbness near the ankle or lateral foot may reflect saphenous or sural nerve irritation and can be temporary or persistent. Recurrent veins may reflect residual tributaries, recanalization, progression of venous disease, or previously unrecognized pelvic or deep reflux. Duplex follow-up is used selectively to evaluate persistent symptoms, suspected thrombosis, or recurrence rather than as a substitute for clinical review.

\section*{References}
\begin{enumerate}
  \item National Institute for Health and Care Excellence. Varicose Veins: Diagnosis and Management. Guideline CG168.
  \item Society for Vascular Surgery and American Venous Forum. Clinical Practice Guidelines for the Management of Varicose Veins.
  \item European Society for Vascular Surgery. Clinical Practice Guidelines on the Management of Chronic Venous Disease of the Lower Limbs.
\end{enumerate}

\clearpage
